\section*{第 \ref{ch:b2:magmatic-series} 章——岩浆系列与构造环境}

\begin{solution}{exo:b2:magmatic-series:1}
（a）$\ce{SiO2} = 49.60$，碱含量 $2.22 + 0.51 = 2.73$：介于 45 与 52 之间，且小于 5，
\emph{玄武岩}（B）。（b）58.5 与 $3.9 + 1.6 = 5.5$：介于 57 与 63 之间；该处上边界
从 (57, 5.9) 到 (63, 7.0) 的线段，其值为 $5.9 + 1.1\times1.5/6 = 6.18$，
在 58.5 处，5.5 位于其下方：\emph{安山岩}（O2）。（c）72.0 与 $4.0 + 4.2 = 8.2$：
超出 69 处竖线，位于 8 以上，且在从 (69, 8) 到 (77, 0) 的直线上方，该直线
在 72 处为 5.0，故为\emph{流纹岩}（R）。（d）47.0 与 $4.5 + 2.5 = 7.0$：在 47 处，区 U1
位于从 (45, 5) 到 (49.4, 7.3) 的直线（其值为 $5 + 2.3\times2/4.4 = 6.05$）与
从 (45, 9.4) 到 (49.4, 7.3) 的直线（其值为 $9.4 - 2.1\times2/4.4 = 8.45$）之间。该岩石
位于 U1 区：标准矿物中橄榄石含量低于 \qty{10}{\percent} 者为\emph{碧玄岩}，
否则（较多时）为\emph{碱玄岩}。
\end{solution}

\begin{solution}{exo:b2:magmatic-series:2}
（a）氧化物之和为 98.40，连同挥发分的总和为 100.00；
\ce{H2O} $1.40 < 2$ 与 \ce{CO2} $0.20 < 0.5$：分析可接受。（b）换算系数
$100/98.40 = 1.01626$: \ce{SiO2} 52.44, \ce{TiO2} 1.22, \ce{Al2O3} 16.36,
\ce{FeO_T} 9.96, \ce{MgO} 6.71, \ce{CaO} 8.33, \ce{Na2O} 3.56, \ce{K2O} 1.42
（和为 100.00）；碱含量 $3.56 + 1.42 = 4.98$。（c）之前：51.60 与 4.90，恰在
52 处玄武岩–玄武安山岩界线之下，为\emph{玄武岩}。之后：52.44 与 4.98；在
52.44 处 O1 的上边界，即从 (52, 5) 到 (57, 5.9) 的线段，其值为 $5 + 0.9\times0.44/5
= 5.08$：\emph{玄武安山岩}（O1）。重算使岩石越过
了边界：靠近界线的岩石，其名称取决于它。（d）在 52.44 处，低钾线为 $0.3 +
0.045\times4.44 = 0.50$，高钾线为 $1.2 + 0.095\times4.44 = 1.62$；1.42
位于二者之间：\emph{中钾}。
\end{solution}

\begin{solution}{exo:b2:magmatic-series:3}
（a）$Q + A + P = 90$：$Q = 25/90 = \qty{27.8}{\percent}$，$100\,P/(A + P) =
100\times35/65 = 53.8$：位于 3b 区，为\emph{二长花岗岩}，属花岗岩。黑云母
与角闪石不计入 QAPF 之和。（b）$Q + A + P = 65$：$Q = 5/65 =
\qty{7.7}{\percent}$（介于 5 与 20 之间）且 $100\times55/60 = 91.7$（高于 90）：区
10*，石英闪长岩、石英辉长岩或石英斜长岩。仅凭 \omterm{def:b2:magmatic-series:qapf}{QAPF 图解}无法
在其中选择；斜长石与镁铁质矿物（此处
\qty{35}{\percent} 为单斜辉石）可以。
\end{solution}

\begin{solution}{exo:b2:magmatic-series:4}
低钾线为 $0.3 + 0.045(\ce{SiO2} - 48)$，高钾线为 $1.2 +
0.095(\ce{SiO2} - 48)$。（a）当 50 时：0.39 与 1.39；$0.40$ 恰在低钾线之上：
\emph{中钾}。（b）当 56 时：0.66 与 1.96：\emph{中钾}。（c）当 62 时：0.93 与 2.53；
$2.90 > 2.53$：\emph{高钾}。（d）当 66 时：1.11 与 2.91；$1.00 < 1.11$：\emph{低钾}。
\end{solution}

\begin{solution}{exo:b2:magmatic-series:5}
（a）玄武岩：$A + F + M = 3.2 + 14.5 + 4.5 = 22.2$；$A = 14.4$，$F = 65.3$，$M = 20.3$。
安山岩：$5.05 + 6.1 + 3.8 = 14.95$；$A = 33.8$，$F = 40.8$，$M = 25.4$。（b）该线
从未超过 $F = 54.5$；玄武岩的 $F = 65.3$，位于 F 一侧
无论 $A$ 取何值：\emph{拉斑质}，属芬纳趋势的富铁玄武岩。对安山岩
在 $A = 33.8$ 处，该线介于 (49.6, 35.0) 与 (52.5, 30.0) 之间，取 $F =
49.6 + 2.9\times1.2/5 = 50.3$；$40.8 < 50.3$：\emph{钙碱性}。（c）磁铁矿，
它在早期即从氧化\omterm{def:b2:magma-emplacement:magma}{岩浆}中带走三价铁，使 $F/M$ 不致
像拉斑玄武岩中那样迅速上升。
\end{solution}

\begin{solution}{exo:b2:magmatic-series:6}
MORB：$\mathrm{Nb_N} = 4/0.658 = 6.08$，$\mathrm{La_N} = 4.12/0.648 = 6.36$，比值
0.956。\omterm{def:b2:elements-in-the-earth:morb-oib}{洋岛玄武岩}：$38/0.658 = 57.8$，$30.8/0.648 = 47.5$，比值 1.215。岛弧
玄武岩：$3/0.658 = 4.56$，$6.1/0.648 = 9.41$，比值 0.484。仅岛弧玄武岩具有
铌异常，即蛛网图上 Nb 处的低谷：铌相对
其邻元素 La 而言亏损：岛弧\omterm{def:b2:magma-emplacement:magma}{岩浆}从板片获得了元素，但
铌很少。MORB 近于平坦；\omterm{def:b2:elements-in-the-earth:morb-oib}{洋岛玄武岩}
的 Nb 多于 La。
\end{solution}

\begin{solution}{exo:b2:magmatic-series:7}
（a）$\mathrm{Ti} = 1.6\times0.5993\times10^4 = \qty{9589}{ppm}$；Ti/100 = 95.9，Zr
110，$3\mathrm Y = 96$；和为 301.9；坐标为 31.8、36.4、31.8：落在 B 区。（b）Ti
16\,780，故得 167.8、220、87；和为 474.8；35.3、46.3、18.3：落在 D 区。（c）Ti 4794，故
得 47.9、50、60；和为 157.9；30.3、31.7、38.0：落在 B 区。玄武岩（b）为板内玄武岩
（属洋岛或大陆）；（a）和（c）位于 B 区，可能为 MORB、岛弧
拉斑玄武岩或钙碱性玄武岩：该图无法判定，还需另一图解
（如蛛网图及 Nb 异常）。
\end{solution}

\begin{solution}{exo:b2:magmatic-series:8}
（a）MORB $142/164 = 0.866$，\omterm{def:b2:elements-in-the-earth:morb-oib}{洋岛玄武岩} $4/173 = 0.023$，岛弧玄武岩 $101/201 =
0.502$。（b）先验概率相等时，由贝叶斯公式得 $0.866/(0.866 + 0.023 + 0.502) =
0.866/1.391 = 0.623$。（c）该玄武岩极不可能为\omterm{def:b2:elements-in-the-earth:morb-oib}{洋岛玄武岩}
（概率 $0.023/1.391 = 0.017$），但无法判定其为 MORB 还是岛弧玄武岩：
大致为五分之三对五分之二。
\end{solution}

\begin{solution}{exo:b2:magmatic-series:9}
每 \qty{100}{g} 的摩尔数：\ce{Na2O} $5.5/61.98 = 0.08874$，\ce{K2O} $4.5/94.20 =
0.04777$，\ce{Al2O3} $10.5/101.96 = 0.10298$。（a）$(0.08874 + 0.04777)/0.10298 =
0.13651/0.10298 = 1.326 > 1$：该流纹岩为\emph{过碱性}。（b）长石可
吸收 $0.10298$ 摩尔的碱性氧化物（共 $0.13651$ 摩尔）；超出部分 $0.03353$，
$0.03353/0.13651 = 0.246$，约占碱质的四分之一，须进入其他矿物。
\end{solution}

\begin{solution}{exo:b2:magmatic-series:10}
（a）$s = (\lambda_F - \lambda_M)/(\lambda_F - \lambda_A)$：对 $\lambda_F = 0.5$，$(0.5
- 0.8)/(0.5 - 0.4) = -3.0$；对 $0.7$，$(-0.1)/(0.3) = -0.333$。（b）由
$s(\lambda_F - 0.4) = \lambda_F - 0.8$ 解得 $\lambda_F = (0.4s - 0.8)/(s - 1)$：
冰岛，$s = -1.10$：$(-0.44 - 0.8)/(-2.10) = 0.590$；喀斯喀特，$s = -0.72$：
$(-0.288 - 0.8)/(-1.72) = 0.633$。（c）喀斯喀特\omterm{def:b2:magma-emplacement:magma}{岩浆}，
$\lambda_F$ 较大，失铁更快：磁铁矿将其带出。（d）$\ln(F/M)$ 以
速率 $\lambda_M - \lambda_F$ 增长，$\ln(A/F)$ 以速率 $\lambda_F - \lambda_A$ 增长；因此
$\ln(A/F)$ 每增加一个单位，$\ln(F/M)$ 即增加 $-s$：冰岛为 1.10，
喀斯喀特为 0.72。$F/M$ 沿两条趋势均上升，但在拉斑系列中升得更快，
即芬纳趋势。
\end{solution}

\begin{solution}{exo:b2:magmatic-series:11}
（a）BHVO-2，每 \qty{100}{g} 的摩尔数：\ce{SiO2} 0.82557，\ce{Al2O3} 0.13182，\ce{FeO}
0.15521, \ce{MgO} 0.18007, \ce{CaO} 0.20328, \ce{Na2O} 0.03580, \ce{K2O} 0.00545.
正长石 0.00545 与钠长石 0.03580；\ce{Al2O3} 余下 $0.13182 - 0.04125 = 0.09057$，
全部与等量的 \ce{CaO} 结合成钙长石；\ce{CaO} 余下 $0.11271$ 进入透辉石；
$(\ce{Mg,Fe})\ce{O}$ 余下 $0.33528 - 0.11271 = 0.22257$ 进入紫苏辉石。所需二氧化硅
量：$6\times0.00545 + 6\times0.03580 + 2\times0.09057 + 2\times0.11271 + 0.22257 =
0.87663$；可用量为 0.82557，缺 0.05106。每个 $(\ce{Mg,Fe})\ce{O}$ 由紫苏辉石转为
橄榄石可节省半个二氧化硅，至多 $0.22257/2 = 0.11129$；转化
$0.10212$ 即已足够，$0.12045$ 仍留在紫苏辉石中。标准矿物中有橄榄石
和紫苏辉石，既无石英也无霞石：岩石
在定义上为\omterm{def:b2:magmatic-series:saturation}{硅饱和}，属含紫苏辉石标准矿物的\emph{橄榄拉斑玄武岩}，
属\omterm{def:b2:magmatic-series:series}{拉斑玄武岩系列}，如夏威夷火山盾状阶段的喷发。（b）摩尔数：0.75732、0.14712、0.16008、0.21092、0.17832、0.06131，
0.01592。或者 0.01592、ab 0.06131、an 0.06988、di 0.10843、hy $0.37100 - 0.10843 =
0.26257$；所需二氧化硅量 $0.09552 + 0.36786 + 0.13976 + 0.21686 + 0.26257 = 1.08257$；
可用量 0.75732，缺 0.32525。将全部紫苏辉石转为橄榄石，可节省
0.13129，余下 0.19396；每摩尔 \ce{Na2O} 由钠长石转为霞石
释放 4 个 \ce{SiO2}，故 $0.19396/4 = 0.04849$
的 0.06131 摩尔钠长石中的 \ce{Na2O} 转为霞石。标准矿物中有霞石：
岩石为硅不饱和，即\emph{碱性玄武岩}，属\omterm{def:b2:magmatic-series:series}{碱性系列}。
\end{solution}

\begin{solution}{exo:b2:magmatic-series:12}
(a) 弧玄武岩 $116/141 = 0.823$，MORB $21/101 = 0.208$，\omterm{def:b2:elements-in-the-earth:morb-oib}{洋岛玄武岩}
$3/137 = 0.022$. (b) $0.823/(0.823 + 0.208 + 0.022) = 0.823/1.053 = 0.782$.
(c) 分子 $0.1\times0.823 = 0.0823$；分母 $0.0823 + 0.6\times0.208 +
0.3\times0.022 = 0.0823 + 0.1248 + 0.0066 = 0.2137$；概率 $0.385$。(d)
对同一岩石做同一检验，所得结果为 0.78 或 0.39，
取决于岩石出露处各构造环境的常见程度：\omterm{def:b2:magmatic-series:discrimination}{判别图解}给出的是似然，
其余则须由地质背景提供。
\end{solution}

\begin{solution}{pb:b2:magmatic-series:1}
\noindent\textbf{第一部分。}

\textbf{1.} 氧化物：熔岩 1 为 99.17，熔岩 2 为 98.65，熔岩 3 为 98.20；计入挥发分后
分别为 99.97、99.80 和 99.45，均接近 100。水（0.70、1.10、1.20）的含量低于
\qty{2}{\percent}，二氧化碳（0.10、0.05、0.05）低于 \qty{0.5}{\percent}：
三者均可接受。

\textbf{2.} 换算系数 $100/99.17 = 1.00837$：\ce{SiO2} 50.82、\ce{Na2O} 2.92、\ce{K2O}
0.63，碱含量 3.55。

\textbf{3.} 熔岩 2：$100/98.65 = 1.01368$，\ce{SiO2} 59.40，\ce{Na2O} 3.95，\ce{K2O}
1.17，碱含量 5.12。熔岩 3：$100/98.20 = 1.01833$，\ce{SiO2} 66.50，\ce{Na2O} 4.48，
\ce{K2O} 1.88，碱含量 6.36。

\textbf{4.} 熔岩 1（50.82，3.55）：玄武岩。熔岩 2（59.40，5.12；该处 O2 区的上界为
$5.9 + 1.1\times2.40/6 = 6.34$）：安山岩。熔岩 3（66.50，6.36；位于 63 与
由 (69, 8) 至 (77, 0) 的连线之间，低于由 (63, 7) 至 (69, 8) 的连线）：英安岩。

\textbf{5.} 原始数据：(50.4, 3.52)、(58.6, 5.05)、(65.3, 6.25)，落在相同的区域。
换算使各点的二氧化硅含量至多改变 \qty{1.2}{\percent}，在此不改变分类，
这与\cref{exo:b2:magmatic-series:2}不同。

\textbf{6.} 依据实测的\omterm{def:b2:magmatic-series:qapf}{实际矿物组成}用 \omterm{def:b2:magmatic-series:qapf}{QAPF 图解}。安山质成分的深成岩
富含斜长石，碱性长石和石英很少，故应落在 10 区或 10*区，
即闪长岩或石英闪长岩（由其\omterm{def:b2:magmatic-series:qapf}{实际矿物组成}决定）。

\noindent\textbf{第二部分。}

\textbf{7.} 熔岩 1：界线值为 0.43 和 1.47，$\ce{K2O} = 0.63$：中钾。熔岩 2：界线值 0.81
和 2.28，1.17：中钾。熔岩 3：界线值 1.13 和 2.96，1.88：中钾。

\textbf{8.} 去挥发分后的 \ce{FeO_T} 和 \ce{MgO}：8.67 与 7.97；6.18 与 3.85；4.07
与 1.63。熔岩 1：$3.55 + 8.67 + 7.97 = 20.19$；$A = 17.6$，$F = 42.9$，$M = 39.5$。
熔岩 2：15.15；33.8、40.8、25.4。熔岩 3：12.06；52.7、33.7、13.5。

\textbf{9.} 熔岩 1，$A = 17.6$：(42.8, 17.1) 与 (46.9, 18.1) 之间的连线在
$F = 42.8 + 4.1\times0.5/1.0 = 44.9$ 处；$42.9 < 44.9$，位于钙碱性一侧，
紧邻该线。熔岩 2，$A = 33.8$：$F = 49.6 + 2.9\times1.2/5 = 50.3$；$40.8 < 50.3$，
属钙碱性。熔岩 3，$A = 52.7$：$F = 40.8 - 2.9\times2.7/5 = 39.2$；$33.7 < 39.2$，
属钙碱性。

\textbf{10.} $\ce{FeO_T}/\ce{MgO}$：$8.67/7.97 = 1.09$，$6.18/3.85 = 1.61$，$4.07/1.63 =
2.50$。比值上升，与所有分异玄武岩相同，但铁本身
从 8.67 降至 4.07，三角图上 $F$ 由 42.9 降至 33.7：无铁富集。
\omterm{def:b2:magmatic-series:series}{拉斑玄武岩系列}则会先向 F 攀升，正如
\cref{exo:b2:magmatic-series:5}中的富铁玄武岩那样。

\textbf{11.} \omterm{def:b2:magmatic-series:series}{钙碱性系列}。磁铁矿自氧化性\omterm{def:b2:magma-emplacement:magma}{岩浆}中结晶，
连同镁一起带走了铁；氧化性物种来自俯冲带的
脱水板片。

\noindent\textbf{第三部分。}

\textbf{12.} $\mathrm{Nb_N} = 2.0/0.658 = 3.04$, $\mathrm{La_N} = 7/0.648 = 10.80$,
$\mathrm{Th_N} = 1.1/0.0795 = 13.84$.

\textbf{13.} $3.04/10.80 = 0.281$，远低于 MORB 的中位数（0.956）和
\omterm{def:b2:elements-in-the-earth:morb-oib}{洋岛玄武岩}的中位数（1.215），甚至低于弧玄武岩的中位数（0.484）：
一个深的铌谷。

\textbf{14.} $13.84/3.04 = 4.55$；蛛网图中 MORB 的中位数为 $2.64/6.08 =
0.43$。与弧玄武岩中富集的其他元素一样，钍是由
俯冲板片加入的。

\textbf{15.} \ce{TiO2} 去挥发分后为 $0.95\times1.00837 = 0.958$；$\mathrm{Ti} =
0.958\times5993 = \qty{5741}{ppm}$；Ti/100 = 57.4，Zr 70，$3\mathrm Y = 54$，求和 181.4：
31.6、38.6、29.8，落在 B 区，为 MORB、弧拉斑玄武岩和钙碱性玄武岩所共有。
它排除了板内玄武岩，但无法在洋脊与
火山弧之间作出判断。该图解是为玄武岩建立的：在安山岩和英安岩中，钛已被氧化物
移除，锆因结晶而富集，它们的区域（熔岩 2 将落在 C 区，熔岩 3 落在所有区域之外）
便毫无意义。

\textbf{16.} 钙碱性、中钾系列，具有深的铌谷和高的
$\mathrm{Th_N}/\mathrm{Nb_N}$，喷发了从玄武岩到英安岩的火山，该火山
形成于俯冲带上方的火山弧中。

\noindent\textbf{第四部分。}

\textbf{17.} 对完全\omterm{def:b2:elements-in-the-earth:partition}{不相容元素}，$C_L = C_0/F$：$F = 0.63/1.17 =
0.538$。

\textbf{18.} 铷 $10/20 = 0.500$；钍 $1.1/2.1 = 0.524$。（\omterm{def:b2:elements-in-the-earth:trace}{微量
元素}按所给分析值计，经换算后
变化均小于 \qty{1}{\percent}。）

\textbf{19.} $C_L/C_0 = F^{D - 1}$：$125/70 = 1.786 = F^{-0.9}$，$F = 1.786^{-1/0.9} =
0.525$。

\textbf{20.} 以钾计，自熔岩 2 至熔岩 3 的 $F = 1.17/1.88 = 0.622$。锶：
$380/480 = 0.792 = 0.622^{D - 1}$，故 $D - 1 = \ln 0.792/\ln 0.622 = -0.234/-0.474 =
0.49$，$D \approx 1.49$：锶已变为\omterm{def:b2:elements-in-the-earth:partition}{相容元素}。斜长石容纳锶，
当时斜长石正大量结晶。

\textbf{21.} 以钾计，$F = 0.63/1.88 = 0.335$（铷 $10/32 = 0.31$，钍
$1.1/3.4 = 0.32$）：约三分之二的玄武岩，$1 - 0.335 = 0.665$，在
英安岩形成之前结晶。

\textbf{22.} 同化的钾会使熔岩 2 中的 \ce{K2O} 高于
仅由结晶所给出的值，故 $F = 0.63/1.17$ 会偏低，而晶体的质量
则偏高。

\textbf{23.} $(0.538 + 0.500 + 0.524 + 0.525)/4 = 2.087/4 = 0.522$；晶体占 $1 - 0.522
= 0.478$ 的玄武岩，即制成安山岩须结晶出 \textbf{\qty[text-series-to-math=true]{48}{kg}
晶体，每 \qty{100}{kg} 玄武岩}。
\end{solution}
